\documentclass{beamer}

\usepackage{default}
\usepackage{xeCJK}
\usepackage{fontspec}
\setmainfont{SimSun}
\usetheme{DarkConsole}

\title{研究报告}
\author{陈泰然}
\institute{USTC nelslip}
\logo{\includegraphics[width=2cm,height=2cm]{figures/ustc_logo_fig}}
%\date{}

\begin{document}
	\frame{\titlepage}
	
%	\begin{frame}
%		\frametitle{目录}
%		\tableofcontents
%	\end{frame}
	
	\AtBeginSection[]
	{
		\begin{frame}
			\frametitle{目录}
			\tableofcontents
		\end{frame}
	}
	
	\section{World Model的核心组成}
	
	\begin{frame}
		\frametitle{World Model的核心组成}
		World Model由三部分组成：
		\begin{itemize}
			\item[V] VAE Model
			\item[M] MDN-RNN Model
			\item[C] Controller Model
		\end{itemize}
	\end{frame}
	
	\begin{frame}
		\frametitle{V Model}
		在World Model中，V的作用是学习输入图像的一个抽象的、高度压缩的表示。在实验中，采用一个简单的Variational Autoencoder（VAE）作为V模型，来将每一帧输入图像编码成为一个很小的向量$z$。
	\end{frame}

	\begin{frame}
		\frametitle{M Model}
		M是一个预测模型，它的作用是预测V模型下一步即将产生的$z$。
	\end{frame}
	
	\begin{frame}
		\frametitle{C Model}
		C的作用是决定agent下一步的动作，让累积回报期望最大。\\
		在实验中，有意地把C做的非常简单，这样做是为了让agent的复杂度集中在V和M中。C是如下的一个线性模型：
		\begin{equation}
			a_t=W_c[z_t\ h_t]+b_c
		\end{equation}
		在该模型中，$W_c$和$b_c$分别是权重矩阵和偏置向量，将$z_t$和$h_t$连接而成的向量映射到动作向量$a_t$。
	\end{frame}

\end{document}
